\nonstopmode{}
\documentclass[letterpaper]{book}
\usepackage[times,inconsolata,hyper]{Rd}
\usepackage{makeidx}
\usepackage[utf8]{inputenc} % @SET ENCODING@
% \usepackage{graphicx} % @USE GRAPHICX@
\makeindex{}
\begin{document}
\chapter*{}
\begin{center}
{\textbf{\huge Package `aFIPC'}}
\par\bigskip{\large \today}
\end{center}
\inputencoding{utf8}
\ifthenelse{\boolean{Rd@use@hyper}}{\hypersetup{pdftitle = {aFIPC: Automated Fixed Item Parameter Linking}}}{}
\ifthenelse{\boolean{Rd@use@hyper}}{\hypersetup{pdfauthor = {Seongho Bae}}}{}
\begin{description}
\raggedright{}
\item[Type]\AsIs{Package}
\item[Title]\AsIs{Automated Fixed Item Parameter Linking}
\item[Version]\AsIs{0.1.0}
\item[Author]\AsIs{Seongho Bae [aut, cre]}
\item[Maintainer]\AsIs{Seongho Bae }\email{seongho@kw.ac.kr}\AsIs{}
\item[Description]\AsIs{Automates fixed item parameter linking for test linking under
the item response theory paradigm using mirt package estimates.}
\item[License]\AsIs{GPL-3 | file LICENSE}
\item[Imports]\AsIs{mirt, methods}
\item[Suggests]\AsIs{testthat (>= 3.0.0)}
\item[Encoding]\AsIs{UTF-8}
\item[Config/testthat/edition]\AsIs{3}
\item[Config/roxygen2/version]\AsIs{8.0.0}
\end{description}
\Rdcontents{\R{} topics documented:}
\inputencoding{utf8}
\HeaderA{autoFIPC}{automated fixed item parameter linking}{autoFIPC}
%
\begin{Description}
automated fixed item parameter linking
\end{Description}
%
\begin{Usage}
\begin{verbatim}
autoFIPC(
  newformXData,
  oldformYData,
  newformCommonItemNames,
  oldformCommonItemNames,
  itemtype = "3PL",
  newformBILOGprior = NULL,
  oldformBILOGprior = NULL,
  tryFitwholeNewItems = T,
  tryFitwholeOldItems = T,
  checkIPD = T,
  tryEM = T,
  freeMEAN = T,
  forceNormalZeroOne = F,
  parameterOverwrite = F,
  empiricalhist = F,
  confirmCommonItems = NULL,
  ...
)
\end{verbatim}
\end{Usage}
%
\begin{Arguments}
\begin{ldescription}
\item[\code{newformXData}] new form data X

\item[\code{oldformYData}] old form (base form) data Y

\item[\code{newformCommonItemNames}] Common item variable names in new form data

\item[\code{oldformCommonItemNames}] Common item variable names in old (base) form data

\item[\code{itemtype}] itemtype of calibration

\item[\code{newformBILOGprior}] using BILOG-MG prior when try to calibrate 3PL model? if you want, set the this to TRUE

\item[\code{oldformBILOGprior}] using BILOG-MG prior when try to calibrate 3PL model? if you want, set the this to TRUE

\item[\code{tryFitwholeNewItems}] do you want to try calibrate new form data? default is TRUE

\item[\code{tryFitwholeOldItems}] do you want to try calibrate base form data? default is TRUE

\item[\code{checkIPD}] do you want to check item parameter drift? default is TRUE

\item[\code{tryEM}] do you want to try EM algorithm when you calibrate model? defalut is TRUE

\item[\code{freeMEAN}] allow freely mean estimation, default is TRUE

\item[\code{forceNormalZeroOne}] set the prior distribution follows N(0,1) distribution. default is FALSE

\item[\code{parameterOverwrite}] don't touch it

\item[\code{empiricalhist}] do you want to use empirical histogram method when tryEM = TRUE? default is FALSE

\item[\code{confirmCommonItems}] set TRUE to accept the supplied common-item pairs without an interactive prompt. Default NULL: in interactive sessions the user will be prompted; set FALSE to explicitly reject (function will stop).

\item[\code{...}] Additional arguments reserved for future extensions.
\end{ldescription}
\end{Arguments}
%
\begin{Value}
the model list of the base form, new form, linked form
\end{Value}
%
\begin{Examples}
\begin{ExampleCode}
## Not run:
autoFIPC(
  newformXData = new_model,
  oldformYData = old_model,
  newformCommonItemNames = common_new,
  oldformCommonItemNames = common_old,
  confirmCommonItems = TRUE
)

## End(Not run)
\end{ExampleCode}
\end{Examples}
\inputencoding{utf8}
\HeaderA{surveyFA}{surveyFA}{surveyFA}
%
\begin{Description}
Fallback calibration helper used when direct \code{autoFIPC()} estimation fails.
\end{Description}
%
\begin{Usage}
\begin{verbatim}
surveyFA(data, autofix = TRUE, forceUIRT = TRUE, forceNormalEM = FALSE,
  forceMHRM = FALSE, unstable = FALSE, SE = TRUE, itemtype = "2PL",
  maxItemRemovals = 3L, pThreshold = 0.05, ...)
\end{verbatim}
\end{Usage}
%
\begin{Arguments}
\begin{ldescription}
\item[\code{data}] A response matrix or data frame.
\item[\code{autofix}] When TRUE, remove least-fitting items and retry calibration.
\item[\code{forceUIRT}] Kept for legacy compatibility; fallback requires TRUE.
\item[\code{forceNormalEM}] Force normal EM/MMLE estimation for fallback attempts.
\item[\code{forceMHRM}] Force MHRM-based estimation for fallback attempts.
\item[\code{unstable}] Use a permissive estimator sequence including QMCEM and MHRM.
\item[\code{SE}] Whether to request standard errors from \code{mirt}.
\item[\code{itemtype}] IRT item type for fallback estimation.
\item[\code{maxItemRemovals}] Maximum number of items to remove in autofix mode.
\item[\code{pThreshold}] p-value cutoff for selecting misfitting items.
\item[\code{...}] Additional arguments reserved for compatibility.
\end{ldescription}
\end{Arguments}
%
\begin{Value}
A fitted \code{mirt} model object with finite log-likelihood and, when
\code{SE = TRUE}, a passing second-order test plus finite covariance estimates.
\end{Value}
\printindex{}
\end{document}
